\section*{Conclusion Générale}

Dans un contexte où le secteur financier marocain connaît une transformation digitale accélérée, la gestion des dossiers de crédit et le recouvrement des impayés nécessitent des outils modernes pour garantir l'efficacité opérationnelle. Les établissements financiers comme EQDOM rencontrent des difficultés pour gérer manuellement un grand volume de dossiers, assurer le suivi des actions de recouvrement et coordonner efficacement les équipes. C'est dans ce contexte que s'inscrit le projet \textbf{Credit Platform EQDOM}, dont la problématique centrale est :

\vspace{0.5cm}

\begin{center}
\textit{Comment concevoir et développer une plateforme web centralisée, sécurisée et intelligente, capable de gérer efficacement l'ensemble du cycle de vie des dossiers de crédit et d'optimiser le processus de recouvrement des impayés, tout en automatisant l'archivage mensuel et en garantissant la traçabilité complète des actions ?}
\end{center}

\vspace{0.5cm}

Pour répondre à cette problématique, nous avons conduit une démarche structurée couvrant la conception, le développement et l'intégration de l'ensemble des composants de la plateforme. La phase de réalisation a abouti à une application web complète, articulée autour d'un backend \textbf{Spring Boot 4.1.0} sécurisé par \textbf{JWT} et d'un frontend \textbf{React 19} avec \textbf{TypeScript}, connectés via une \textbf{API REST} documentée.

La plateforme couvre l'ensemble des besoins identifiés : les \textbf{administrateurs} disposent d'un espace complet pour superviser les dossiers, importer des fichiers Excel, affecter les dossiers aux agents et lancer l'archivage mensuel automatique. Les \textbf{agents} bénéficient d'un tableau de bord personnalisé pour gérer leurs dossiers assignés, enregistrer leurs actions avec traçabilité complète et suivre leurs performances en temps réel.

L'élément innovant majeur du projet réside dans le \textbf{système d'archivage mensuel automatique avec détection intelligente des récurrences}. À la fin de chaque mois, le système archive l'historique complet, génère des exports CSV et sauvegarde les identifiants pour la détection. Lors de l'import suivant, le moteur intelligent détecte automatiquement les clients récurrents et applique les statuts appropriés selon l'historique (\texttt{CONTACTE} $\rightarrow$ \texttt{RELANCE}, \texttt{PROMESSE} $\rightarrow$ \texttt{SUIVI\_PROMESSE}), garantissant ainsi la continuité du suivi client et la conservation d'un historique illimité.

Au terme de ce projet, nous pouvons affirmer que les objectifs initiaux ont été pleinement atteints. \textbf{Credit Platform EQDOM} constitue une réponse concrète aux lacunes des systèmes traditionnels, en proposant une solution globale, automatisée et sécurisée.

Sur le plan technique, ce projet nous a permis de consolider des compétences essentielles : développement \textbf{full-stack} avec \textbf{Spring Boot} et \textbf{React}, sécurisation des API par \textbf{JWT} et \textbf{Spring Security}, modélisation de bases de données avec \textbf{PostgreSQL} et fonctions \textbf{PL/pgSQL}, conception d'une machine d'états rigoureuse, traitement de fichiers Excel avec \textbf{Apache POI}, et déploiement cloud via \textbf{Docker}, \textbf{Railway} et \textbf{Vercel}. Il a renforcé notre capacité à conduire un projet de bout en bout, de l'analyse des besoins jusqu'à la livraison d'une application opérationnelle en production.

Plusieurs axes d'amélioration restent envisageables pour les versions futures :

\begin{itemize}
    \item L'intégration de notifications en temps réel via \textbf{WebSockets}
    \item L'implémentation d'\textbf{intelligence artificielle} pour la prédiction du risque d'impayé
    \item L'ajout d'un système de \textbf{reporting avancé} avec tableaux de bord personnalisables
    \item L'orchestration \textbf{Kubernetes} pour la haute disponibilité
    \item Le développement d'une application mobile \textbf{Progressive Web App (PWA)} pour les agents en déplacement
\end{itemize}

\textbf{Credit Platform EQDOM} représente ainsi pour nous bien plus qu'un projet académique : c'est une démonstration concrète de notre capacité à concevoir et réaliser une solution numérique complète, moderne et utile, ancrée dans une problématique réelle du secteur financier marocain. Plus qu'un outil de gestion, cette plateforme constitue un levier stratégique pour EQDOM dans sa démarche de transformation digitale, contribuant à l'amélioration de l'efficacité du recouvrement, à la réduction des impayés par un suivi rigoureux et automatisé, et à l'optimisation de la productivité des équipes dans un environnement économique en constante évolution.
